% Substitua o texto abaixo pelo seu resumo.



Este trabalho apresenta o projeto e a implementação de uma infraestrutura baseada em Sistemas Embarcados voltada para o monitoramento centralizado de experimentos em laboratórios didáticos de Física. A motivação fundamenta-se nas barreiras estruturais de laboratórios tradicionais, que frequentemente operam de forma isolada e exigem hardware dedicado para cada bancada, dificultando a gestão pedagógica e a escalabilidade. A arquitetura proposta utiliza microcontroladores ESP32 como nós de borda para a aquisição de dados e um servidor Raspberry Pi como coordenador central. A metodologia seguiu uma abordagem de desenvolvimento tecnológico, avaliada por meio de uma Prova de Conceito (PoC) com o experimento de Carga e Descarga de Capacitor (Circuito RC). A comunicação entre os dispositivos ocorre via protocolo HTTP/JSON sobre rede Wi-Fi, eliminando a necessidade de conexões cabeadas e centralizando a telemetria em um banco de dados PostgreSQL. O sistema conta com uma interface web desenvolvida em Python (Quart), permitindo a visualização em tempo real e o download de dados por alunos e professores. Os resultados demonstram a viabilidade de uma camada de orquestração que desacopla a interface de usuário do hardware de coleta, promovendo um ambiente laboratorial integrado, escalável e capaz de suportar modernização da infraestrutura laboratorial.

 \textbf{Palavras-chave}: Sistemas Embarcados; Experimentação Sem Fio; Ensino de Física; ESP32; Telemetria.

